\subsection{连通空间与连通集}

定义 1. 如果拓扑空间 X 不是两个不相交的非空开集的并，就称它是连通的。

把“开集”一词换成“闭集”，便得到一个等价的定义。X 连通当且仅当 X 中既开又闭的子集只有空集与全空间 X 两者。

若 X 连通，且 A, B 是两个非空的开（相应地，闭）子集，使得 \(A \cup B = X\)，则 \(A \cap B \neq \varnothing\)。

例. * 1) 我们将在第四章，§ 2，第 5 目看到，实直线是连通的，而有理直线不是连通的。*

\begin{enumerate}
\def\labelenumi{\arabic{enumi})}
\setcounter{enumi}{1}
\tightlist
\item
  多于一个点的离散空间不是连通的。
\end{enumerate}

注意，若 \((\mathbf{U}_{t})_{t\in\mathbf{I}}\) 是由开集组成的拓扑空间 X 的一个分拆{[}由分拆的定义（见《集合论》，R，§ 4，第 4 目），这些开集必然非空{]}，则每个 \(U_{t}\) 在 X 中既开又闭，因为 \(U_{t}\) 的余集是诸 \(U_{x}\)（其中 \(x\neq t\)）的并。X 的开子集正是这样的子集 A：\(A\cap U_{t}\) 在 \(U_{t}\) 中是开的（对每个 \(t\in I\)），因而 X 可等同于诸 \(U_{t}\) 的和（§ 2，第 4 目，例 3），且当 I 多于一个元素时 X 不连通。

定义 2. 如果拓扑空间 X 的子集 A 作为 X 的子空间是连通的，就称 A 为一个连通集。

A 是 X 的连通子集的必要充分条件是：对 A 的每个由 X 的两个开（或闭）子集 B, C 组成的覆盖，只要 \(A \cap B\) 与 \(A \cap C\) 都非空，就有 \(A \cap B \cap C \neq \emptyset\)。

例. 在任何拓扑空间中，空集与每个由单独一点组成的集合都是连通的。在豪斯多夫空间 X 中，多于一个点的有限集都不连通；更一般地，X 中每个多于一个点且至少有一个孤立点的子集都不连通。

若稠密子集 A 连通，则全空间 X 连通；否则将存在 X 的两个非空不相交的开子集 M, N，使得 \(M \cup N = X\)，而 \(M \cap A\), \(N \cap A\) 将是 A 的两个不相交的非空开子集，其并为 A。于是我们有

命题 1. 若 A 是连通集，则每个满足 \(A \subset B \subset \overline{A}\) 的集合 B 都是连通的。

命题 2. 交非空的连通集族之并是连通的。

设 \((\mathbf{A}_t)_{i\in I}\) 是 \(\mathbf{X}\) 的连通子集的一个族，它们都含有同一个点 \(x\)；我们须证明

\[
\mathbf {A} = \bigcup_ {i} \mathbf {A} _ {i}
\]

是连通的。若不然，存在两个开集 B 与 C，使得 \(B \cap A\) 与 \(C \cap A\) 非空，且 \(A \subset B \cup C\) 及 \(A \cap B \cap C = \varnothing\)。x 属于 B, C 两集之一，设 \(x \in B\)；另一方面，诸集合 \(A_t\) 之一，设为 \(A_x\)，与 C 相交；因此我们有 \(A_x \subset B \cup C\)，\(A_x \cap B \cap C = \varnothing\)，且 \(B \cap A_x\) 与 \(C \cap A_x\) 非空。于是 \(A_x\) 不连通，这就得出矛盾。

推论. 设 \((\mathbf{A}_n)_{n\geqslant 0}\) 是连通集的一个无穷序列，使得 \(\mathbf{A}_{n + 1}\cap \mathbf{A}_n\neq \emptyset\) 对每个 \(n\geqslant 0\) 成立。则并 \(\bigcup_{n = 0}^{\infty}\mathbf{A}_n\) 是连通的。

对 \(n\) 用归纳法，由命题 2 立即看出，集合 \(\mathbf{B}_n = \bigcup_{i=0}^n \mathbf{A}_i\) 对每个 \(n\) 都是连通的。诸集合 \(\mathbf{B}_n\) 的交非空；故其并（等于 \(\bigcup_{n=0}^{\infty} \mathbf{A}_n\)）由命题 2 是连通的。

命题 3. 设 A 是拓扑空间 X 的一个子集。若 B 是 X 的连通子集，且同时与 A 和 \(\mathbb{C}\) A 相交，则 B 与 A 的边界相交。

若不然，B 与 A 的内部及外部的交将是 B 的两个开子集，且构成 B 的一个分拆，从而 B 将不连通。

推论. 在连通空间 X 中，每个异于 X 自身的非空集都至少有一个边界点。

